\chapter{Общий носитель гиперзарядового breaking-background}
\label{ch:v10-h15-r2-hypercharge-background-carrier}

Аудит коэффициентов оставил один ближайший маршрут: единый spectral trace
на фоне, нарушающем Pati--Salam до Стандартной модели. Наивная реализация
двумя независимыми adjoint-полями не минимальна. Общий носитель уже имеется
в constrained composite branch:
\begin{equation}
 \boxed{\Delta\sim(\mathbf2_R,\mathbf1_L,\mathbf4_4)}.
\end{equation}

В базисе трёх цветных компонент и одной лептонной компоненты четвёрки
спектры $6T_R^3$ и $3(B-L)$ на восьми компонентах $\Delta$ равны
\begin{align}
 t_\Delta&=(3,3,3,-3,-3,-3,3,-3),\\
 b_\Delta&=(1,1,1,1,1,1,-3,-3).
\end{align}
Их сумма даёт
\begin{equation}
 q_\Delta=6Y=(4,4,4,-2,-2,-2,0,-6).
\end{equation}
Нейтральное направление единственно. Его точный проектор есть
\begin{equation}
 P_0=-\frac{(Q_\Delta-4I)(Q_\Delta+2I)(Q_\Delta+6I)}{48},
 \qquad \rank P_0=1.
\end{equation}

Вес нейтральной компоненты относительно двух Cartan-генераторов равен
$(3,-3)$. Для генератора $\alpha(6T_R^3)+\beta,3(B-L)$ условие сохранения
VEV имеет вид
\begin{equation}
 3\alpha-3\beta=0.
\end{equation}
Его ядро одномерно и порождается $(1,1)$. Поэтому единственный неснятый
Cartan-луч действует как
\begin{equation}
 \boxed{6Y=6T_R^3+3(B-L)}.
\end{equation}
Один $\Delta$-VEV тем самым заменяет пару независимо подобранных adjoint
spurions и автоматически сохраняет $SU(2)_L$.

На восьми SM-секторах $\Sigma$ два исходных вектора зарядов равны
\begin{align}
 t_\Sigma&=(3,-3,3,-3,3,-3,3,-3),\\
 b_\Sigma&=(0,0,4,4,-4,-4,0,0).
\end{align}
Они независимы и ортогональны, с Gram-матрицей
$\operatorname{diag}(72,64)$. Луч $(1,1)$ восстанавливает прежний спектр
\begin{equation}
 q_\Sigma=(3,-3,7,1,-1,-7,3,-3).
\end{equation}
Алгебра диагональных функций от $(T_R^3,B-L)$ имеет ранг $6$; алгебра
полиномов от $Y$ также имеет ранг $6$, а их объединение ранга не повышает.
Значит, на $\Sigma$ они задают одно и то же спектральное разложение и
действительно содержат ранее найденный $R_2$-проектор и gap.

Статус, однако, условный. В general fundamental scalar menu поля $\Delta$
нет; оно появляется после composite first-order restriction. Минимальный
composite spectral potential делает желаемый rank-one orbit седлом с шестью
отрицательными модами. Ранее построенный irreducible relative mapping-cone
parent исправляет $\Delta+C$-Hessian до сигнатуры $(43_+,9_0,0_-)$, но ещё
не встроен в текущий Callias--$M_4$--$H_{15}$ parent Тома X.

\begin{theorem}[Минимальный общий носитель гиперзаряда]
Нейтральный rank-one луч поля
$\Delta=(2_R,1_L,4_4)$ имеет единственный Cartan-стабилизатор
$Y=T_R^3+(B-L)/2$. Его действие на $\Sigma$ порождает ту же шестимерную
спектральную алгебру, что и совместное действие $T_R^3$ и $B-L$.
Следовательно, отдельные adjoint-фоны для построения гиперзарядового gap не
требуются.
\end{theorem}

\begin{nogocriterion}[Стабилизатор не равен vacuum-parent]
Наличие нейтральной компоненты и правильной неподвижной подгруппы не
доказывает существование устойчивого VEV. До типизированного вложения
relative mapping-cone curvature в общий родитель нельзя переносить его
положительный Hessian или фиксированный коэффициент в ветвь Тома X.
\end{nogocriterion}

\begin{protocol}\begin{enumerate}
\item размерность $\Delta$: \textbf{$8$ complex};
\item нейтральных направлений: \textbf{$1$};
\item stabilizer-kernel: \textbf{$\operatorname{span}(1,1)$};
\item неснятый генератор: \textbf{$6Y=6T_R^3+3(B-L)$};
\item ранг joint/hypercharge spectral algebra: \textbf{$6/6$};
\item общий scalar carrier: \textbf{$\Delta$, условно допущен};
\item general-fundamental/composite presence: \textbf{$0/1$};
\item inherited Tome X background-map rank: \textbf{$0$};
\item conditional/current origin: \textbf{$3/4$ и $2/4$};
\item следующий гейт: \textbf{typed embedding $\Delta$ mapping-cone parent}.
\end{enumerate}\end{protocol}

Машинный сертификат сохранён в
\path{s2t_v10_particle_wrinkle_dislocation_callias_equal_charge_twist_m4_h15_r2_hypercharge_breaking_background_common_carrier_admission_gate_results.json}.